\section{错误模型的代价：从交叉熵到 KL 散度}
\label{sec:07-Information-Theory/03-Divergences/01}

你拿到了一个压缩任务。消息只有三种：\(a,b,c\)。有人告诉你它们的真实出现频率是

\[
P = (1/2,\; 1/4,\; 1/4).
\]

如果你知道 \(P\)，可以给 \(a\) 分配短码、给 \(b\) 和 \(c\) 分配稍长的码。理想信息量 \(-\log_2 P(x)\) 分别是 \(1, 2, 2\) bit，平均码长

\[
\frac12 \cdot 1 + \frac14 \cdot 2 + \frac14 \cdot 2 = 1.5 \text{ bit}.
\]

现在换一个编码者。他不了解真实分布，凭直觉假设三种消息等可能：

\[
Q = (1/3,\; 1/3,\; 1/3).
\]

他按 \(-\log_2 Q(x) = \log_2 3 \approx 1.585\) bit 给每个符号分配码长。因为消息 \(a\) 实际出现概率是 \(1/2\) 而不是 \(1/3\)，码长分配偏离了最优。平均下来，每个符号他需要花费

\[
\frac12 \cdot 1.585 + \frac14 \cdot 1.585 + \frac14 \cdot 1.585 = 1.585 \text{ bit}.
\]

比最优方案多了约 \(0.085\) bit。

这个差距看起来很小，但它承载了一个普遍问题：当模型与数据不一致时，我们的编码、预测、决策都会付出额外代价。接下来的任务是把这种额外代价精确地分离出来——算清其中多少来自数据本身的不可压缩性，多少来自模型的错误。

\subsection{交叉熵：用 Q 的眼光度量 P 的数据}

如果用分布 \(Q\) 为来自分布 \(P\) 的数据分配理想码长，平均码长是

\[
H(P, Q) = \E_{X\sim P}[-\log Q(X)] = -\sum_x P(x) \log Q(x).
\]

这个量叫交叉熵。注意两个分布的角色完全不同：\(P\) 在求和外面当权重，决定哪些结果频繁出现；\(Q\) 在对数里面，决定每个结果被分配多长的码。把 \(P\) 和 \(Q\) 的位置互换，算的是另一个问题——用 \(P\) 的眼光去度量 \(Q\) 的数据。

同一个公式直接就是机器学习里负对数似然的原型。模型给观测 \(x_i\) 的概率是 \(q_\theta(x_i)\)，训练集上的平均损失

\[
-\frac{1}{n}\sum_{i=1}^n \log q_\theta(x_i)
\]

恰好是在用经验分布近似 \(P\) 时估计交叉熵。最大化似然和最小化交叉熵是同一个动作的两种叫法。

\subsubsection{分类任务里为什么只剩一项}

在 \(K\) 类分类中，一个样本的真实类别是 \(y\)，对应的 one-hot 分布满足 \(P(y)=1\)，其余类别概率为零。模型输出预测概率 \(q_1,\ldots,q_K\)，代入交叉熵：

\[
H(P, Q) = -\sum_{k=1}^K P(k) \log q_k = -\log q_y.
\]

因为只有真实类别的权重非零，其他项全部消失。但这不意味着模型对其他类别的预测无关紧要——\(q_y\) 的值受 softmax 分母中所有类别得分的影响，每个预测都参与了归一化竞争。

注意负对数对``自信但错误''的惩罚极重。以自然对数计，把真实类别的预测概率从 \(0.9\) 压到 \(0.1\)，损失从约 \(0.105\) 飙到约 \(2.303\)——增加了二十倍。这不是线性惩罚，是指数级的厌恶。

\subsection{减去数据本身的熵，剩下的就是模型要付的账单}

交叉熵里至少有一部分不是模型的错——数据本身就有不可消除的不确定性。从交叉熵中减去熵，把数据固有的部分剥离：

\[
\begin{aligned}
H(P, Q) - H(P)
&= -\sum_x P(x) \log Q(x) + \sum_x P(x) \log P(x) \\
&= \sum_x P(x) \log \frac{P(x)}{Q(x)}.
\end{aligned}
\]

这个差值叫 Kullback--Leibler 散度，或相对熵：

\[
\KL(P \| Q) = \sum_x P(x) \log \frac{P(x)}{Q(x)}.
\]

于是最重要的分解出现了：

\[
H(P, Q) = H(P) + \KL(P \| Q).
\]

回到开头那个三消息的例子。\(H(P) = 1.5\) bit，\(H(P, Q) \approx 1.585\) bit，所以

\[
\KL(P \| Q) \approx 0.085 \text{ bit}.
\]

这个 \(0.085\) bit 的含义非常具体：长期使用模型 \(Q\) 为来自 \(P\) 的数据编码，在长序列算术编码的渐近极限下，每个符号比知道真实分布多花约 \(0.085\) bit。它不是抽象距离，是多付的账单。

\subsection{一个模型说某事概率为零——这不是小错}

假设真实分布里某个结果确实会出现：\(P(x) > 0\)。但模型断言它不可能：\(Q(x) = 0\)。此时

\[
-\log Q(x) = +\infty,
\]

交叉熵和 KL 散度都是无穷。这不是数值计算上的溢出问题——概率模型中的零是一条硬性断言。它不是在说``这件事非常罕见''，而是在说``这件事绝无可能''。只要现实违背一次，整个编码方案就没有为该结果预留任何有限表示。

反过来，若 \(P(x) = 0\) 而 \(Q(x) > 0\)：这个结果从不由 \(P\) 产生，所以它对 \(P\) 下的平均代价没有贡献。符号上通常约定 \(0 \log(0/q) = 0\)；但 \(p \log(p/0) = +\infty\) 对任何 \(p > 0\) 都不会改变。

用测度论的语言重述一遍：KL 有限要求 \(P\) 对 \(Q\) 绝对连续，记作 \(P \ll Q\)——凡是 \(Q\) 判定概率为零的事件，\(P\) 也必须判为零。这条支持条件不是技术细节；它决定了在什么情况下 KL 散度可以作为一个有限工具使用。

\subsection{不对称不是缺陷}

取最极端的离散例子：

\[
P = (1, 0), \qquad Q = (1/2, 1/2).
\]

以 2 为底：

\[
\KL(P \| Q) = 1 \cdot \log_2 \frac{1}{1/2} + 0 \cdot \log_2 \frac{0}{1/2} = 1 \text{ bit}.
\]

\(Q\) 把一半的码空间浪费在了根本不会出现的第二个结果上——多花了 1 bit。

现在把方向反过来：

\[
\KL(Q \| P) = \frac12 \log_2 \frac{1/2}{1} + \frac12 \log_2 \frac{1/2}{0} = +\infty.
\]

因为 \(Q\) 有时会产生第二个结果，而 \(P\) 对它分配了零概率。同一个分布对，两个方向的 KL 一个有限、一个无穷。

这不是 KL 的``缺陷''。两个方向回答的是两个不同的问题：用宽松模型去拟合集中数据，与用过度集中的模型去描述较分散的数据，面临的惩罚结构完全不同。编码的视角把这种不对称翻译成了实在的代价，不对称性反映了任务本身的不对称。

因此 \(D(P\|Q)\) 不满足对称性，通常也不满足三角不等式。称它为``散度''比称它为``距离''更准确。

\subsection{连续情形的 KL 为什么比微分熵更稳}

若 \(P, Q\) 有密度 \(p, q\)，定义照搬：

\[
\KL(P \| Q) = \int p(x) \log \frac{p(x)}{q(x)} \, dx.
\]

前面学过，单独的微分熵会随坐标缩放而改变——如果你把 \(x\) 换成 \(2x\)，\(h(X)\) 的值会变。但在 KL 中，换元时 \(p\) 和 \(q\) 同时乘上同一个 Jacobian 因子，密度比中的坐标因子相消：

\[
\frac{p(u)}{q(u)} = \frac{p(x)|\det J|^{-1}}{q(x)|\det J|^{-1}} = \frac{p(x)}{q(x)}.
\]

所以 KL 在可逆光滑换元下保持不变。这是连续场景里相对量比单个微分熵更可靠的原因之一。

测度层面的统一写法是

\[
\KL(P \| Q) = \E_P\!\left[\log \frac{dP}{dQ}\right],
\]

其中 \(dP/dQ\) 是 Radon--Nikodym 导数。这个表达式再次提醒：期望永远在 \(P\) 下取——是第一参数决定了对哪些区域加权。

\subsection{一个尚未证明但编码直觉强烈担保的事实}

回到开头：知道真实分布后设计的最优码，怎么可能被一个错误模型打败？如果KL散度可以为负，那就意味着用错误的 \(Q\) 编码，有时反而比用真实的 \(P\) 编码更省 bit——这违背了``真实分布知道数据长什么样''的基本直觉。

但光有直觉不够。KL 求和式里，当某个 \(x\) 上 \(Q(x) > P(x)\) 时，\(\log[P(x)/Q(x)]\) 是负的。这些负项被 \(P(x)\) 加权后，为什么总和不可能为负？这个事实需要一个不等式来锁定。下一节做这件事。
